\documentclass[10pt,letterpaper]{article}

% Basic packages
\usepackage[margin=1in]{geometry}
\usepackage[utf8]{inputenc}
\usepackage{hyperref}

% Hyperref setup - clean blue links
\usepackage{xcolor}

% define a softer elegant blue (adjust RGB or HTML code as you like)
\definecolor{ElegantBlue}{HTML}{1A4E8A} % deep, refined blue

\hypersetup{
    colorlinks=true,
    linkcolor=ElegantBlue,
    urlcolor=ElegantBlue,
    citecolor=ElegantBlue,
    pdftitle={Andrea Ciccarone - CV},
    pdfauthor={Andrea Ciccarone}
}


% Remove page numbers
\pagestyle{empty}
\usepackage{tabularx}

% Adjust spacing
\usepackage{setspace}
\setlength{\parindent}{0pt}
\setlength{\parskip}{0pt}

% Section formatting with better spacing
\usepackage{titlesec}
\titleformat{\section}{\normalfont\large\bfseries}{\thesection}{0pt}{}[{\titlerule[0.5pt]}]
\titlespacing{\section}{0pt}{12pt}{8pt}

\titleformat{\subsection}{\normalfont\normalsize\bfseries}{\thesubsection}{0pt}{}
\titlespacing{\subsection}{0pt}{8pt}{5pt}

\begin{document}

% Header
\begin{center}
{\Large \textbf{Andrea Ciccarone}}\\[3pt]
Columbia Business School, Columbia University\\
665 W 130th St, New York, NY 10027\\
\href{mailto:ac4790@columbia.edu}{ac4790@columbia.edu} \quad (646) 525-0556\\
\href{https://andreaciccarone.github.io}{andreaciccarone.github.io}
\end{center}

\vspace{5pt}

\begin{center}
\begin{tabular}{p{0.46\linewidth} p{0.46\linewidth}}
\textbf{Placement Chairs:} & \textbf{Placement Administrators:} \\[3pt]
Mark Dean (\href{mailto:mark.dean@columbia.edu}{mark.dean@columbia.edu}) &
Jonathan Mendoza (\href{mailto:jam2546@columbia.edu}{jam2546@columbia.edu}) \\[3pt]
Martin Uribe (\href{mailto:mu2166@columbia.edu}{mu2166@columbia.edu}) &
Amy Devine-Keum (\href{mailto:aed2152@columbia.edu}{aed2152@columbia.edu}) \\
\end{tabular}
\end{center}

\vspace{3pt}


\section*{Education}

\textbf{Columbia University} \hfill 2020 - 2026\\
Ph.D. Candidate in Economics (BusEc track)

\medskip

\textbf{Bocconi University} \hfill 2020\\
M.S. in Economics and Social Sciences, \textit{110/110 Cum Laude}

\medskip

\textbf{Bocconi University} \hfill 2017\\
B.S. in International Economics and Finance, \textit{110/110 Cum Laude}

\section*{Research Interests}

Political Economy, Media Economics, Organizational Economics

\section*{Job Market Paper}

\textbf{An Image is Worth 1000 Words? Partisanship and Attention in Videos}

\medskip

\textit{Abstract:} Political news today is increasingly consumed in attention-scarce, video-based environments. Yet existing research on media partisanship remains text-centric and implicitly assumes long-form, attentive processing. This paper studies how partisanship operates in video news through the transmission of multimodal signals—images and text—under limited attention. I develop a multimodal measure of video partisanship to quantify partisan content in video news and decompose the contribution of each modality. I show that the informational strength of each modality depends on attention: images convey partisanship rapidly through affective cues, while text operates through substantive information that requires sustained exposure to accumulate. A survey experiment using real news footage shows how these properties shape viewers' responses to political videos. Partisan images elicit immediate emotional responses even under brief exposure, whereas partisan text shifts policy attitudes only with sustained exposure. Overall, the results show how low-attention video environments reshape persuasion, by filtering out substantive textual information and amplifying affective visual signals. Efforts to improve media quality must account for the asymmetric roles of visual and textual content.

\section*{Working Papers}

\textbf{Fixed Effects Topic Model}\\
\textit{with \href{https://dan-biderman.netlify.app/}{Dan Biderman}, \href{https://sites.google.com/view/raissafabregas}{David Blei}, \href{https://sites.google.com/view/raissafabregas}{Wei Cai}, \href{https://sites.google.com/view/raissafabregas}{Amir Feder} and \href{https://sites.google.com/view/raissafabregas}{Andrea Prat}}

\medskip

\textit{Abstract:} Social scientists wish to perform topic modeling on documents that are created by different authors in different contexts. However, the same broad topic may be expressed in different ways depending on the environment where the author operates. For example, one may wish to use employee reviews to identify broad corporate culture topics, but the language of reviews is influenced by industry-specific jargon. Existing methods attempt to control for these biases ex post, such as with traditional fixed effect regressions. But these methods cannot fully separate global themes from category-specific language within them. In this paper, we introduce the Fixed Effects Topic Model (FETM), a novel approach to disentangling broad
topics from contextual influences by incorporating
fixed effects directly into the generative process of language. We use
the FETM to identify themes in a large corpus of Glassdoor job reviews. We show
that it outperforms conventional topic models, both in
interpretability and predictive accuracy.

\vspace{10pt}

\textbf{Bonus or Budget? Equity-Efficiency Tradeoff in Performance-Based Transfers}\\
\textit{with \href{https://www.luigicaloi.com/}{Luigi Caloi}}

\medskip

\textit{Abstract:} To improve service delivery, central governments can tie intergovernmental transfers to local policy performance. While performance-based transfers incentivizes local governments and generate efficiency gains, they also shift transfers from low- to high-capacity governments, creating equity losses. We study this equity-efficiency trade-off using a bundle of transfer reforms to Brazilian municipalities. When two states tied transfers to relative educational performance, student test scores rose substantially: moving from the 25th to the 75th percentile of per capita conditional transfers increased scores by 0.16 standard deviations. However, the reform also widened funding disparities, which translated to disparities in expenditures across sectors. In contrast, contemporaneous reforms to unconditional transfers had negligible effects on student outcomes. We use a simple model of optimal transfers to interpret these findings. Our results suggest that the introduction of performance-based transfers delivered large efficiency gains, limited equity costs, and was welfare enhancing. We find minimal evidence of multitasking distortions or score manipulation. Instead, we document improvements in the quality of education inputs and suggestive evidence of reduced corruption.


\section*{Awards, Fellowships, and Grants}

CPE Graduate Student Grant Opportunity, \textit{Columbia Center for Political Economy}  \hfill 2024

\smallskip

CELSS Seed Grant, \textit{Columbia Experimental Laboratory in the Social Sciences}  \hfill 2024

\smallskip

J-PAL's Jobs and Opportunity Initiative Brazil (with Luigi Caloi), \textit{J-PAL} \hfill 2023

\smallskip

Wueller Fellowship for Best Teaching Fellow (Runner-up), \textit{Columbia University} \hfill 2022

\smallskip

Dean's Fellowship, \textit{Columbia University} \hfill 2020 – Present

\section*{Seminars \& Conferences}
IOG-BFI Spring Conference, Economics of Media Bias Workshop, \hfill 2026\\
Bocconi AI and Society Conference, CEPR Media Workshop

\smallskip

AI+Economics Workshop Zurich \hfill 2025


\section*{Other Research Experience}

Research Assistant to Prof. Andrea Prat and Andrey Simonov, \textit{Columbia Business School} \hfill 2021--2024

\smallskip

Research Assistant to Prof. Cailin Slattery, \textit{Columbia Business School} \hfill 2021--2022

\smallskip

Research Assistant to Prof. Thomas Le Barbanchon, Bocconi University \hfill 2019--2020

\smallskip

Research Assistant to Prof. Marco Ottaviani, Bocconi University \hfill 2019--2020

\section*{Teaching Experience}

Main Instructor in Math Camp (Economics MA) \hfill 2022 -- 2024

\smallskip

Guest Lecturer in Political Economy (PhD), Prof. Andrea Prat \hfill 2022 -- 2024

\smallskip

TA in Industrial Organization (MA), Prof. Ildikó Magyari \hfill 2024

\smallskip

TA in Game Theory and Business (MBA), Prof. Andrea Prat \hfill 2022 -- 2024

\smallskip

TA in Corporate Finance (UG), Prof. Tri Vi Dang \hfill 2023

\smallskip

TA in Math Methods for Economists (MA), Prof. Isaac Bjorke \hfill 2022

\section*{Academic Service}

Referee: \textit{Economica, Economic Journal }

\section*{Bio}

Citizenship: Italian \\
Gender: Male

\section*{References}

% tighten columns a bit (optional)
\setlength{\tabcolsep}{6pt}

\begin{tabular}{@{}p{0.33\linewidth}p{0.33\linewidth}p{0.33\linewidth}@{}}
\raggedright\textbf{Andrea Prat}\\
Professor of Economics \& Business\\
Columbia University\\
\href{mailto:andrea.prat@columbia.edu}{andrea.prat@columbia.edu}
&
\raggedright\textbf{Suresh Naidu}\\
Professor of Economics \& International and Public Affairs\\
Columbia University\\
\href{mailto:sn2430@columbia.edu}{sn2430@columbia.edu}
&
\raggedright\textbf{Mark Dean}\\
Professor of Economics\\
Columbia University\\
\href{mailto:mark.dean@columbia.edu}{mark.dean@columbia.edu}
\\
\end{tabular}






\end{document}
